%% Supplementary Information -- v2 (TSEP)
\documentclass[11pt]{article}
\usepackage[a4paper,margin=2.3cm]{geometry}
\usepackage[T1]{fontenc}\usepackage{lmodern}\usepackage{microtype}
\usepackage{amsmath,amssymb,booktabs,array,tabularx,caption,float,graphicx,xcolor}
\captionsetup{font=small,labelfont=bf,labelsep=period}
\usepackage[numbers,sort&compress]{natbib}\bibliographystyle{unsrtnat}
\usepackage{pdflscape}
\usepackage{hyperref}\hypersetup{colorlinks=true,linkcolor=black,citecolor=blue,urlcolor=blue}
\renewcommand{\thetable}{S\arabic{table}}\renewcommand{\thefigure}{S\arabic{figure}}
\renewcommand{\theequation}{S\arabic{equation}}\renewcommand{\thesection}{S\arabic{section}}
\newcommand{\TODO}[1]{\textcolor{red}{[#1]}}
\newcommand{\LNtwo}{LN$_2$}\renewcommand{\Re}{\mathit{Re}}\newcommand{\Nu}{\mathit{Nu}}
\graphicspath{{../figures/}}
\title{\textbf{Supplementary Information}\\[4pt]\large Helium gas or liquid nitrogen for the first thermal
shield of a 4\,K cryostat: conjugate heat-transfer budget, coolant operating maps and work ledger}
\author{Imerson Joao}\date{}
\begin{document}\maketitle

\noindent All cases, scripts and tabulated results: \href{https://github.com/Imerson/cryostat-first-shield-He-vs-LN2}{github.com/Imerson/cryostat-first-shield-He-vs-LN2}.

\section{Enclosure model}

\begin{table}[H]\centering\small
\caption{Regions, patches and boundary conditions of the enclosure CHT cases.}
\begin{tabularx}{\textwidth}{l l l X}
\toprule
Region & Patch & $T$ condition & Radiation condition \\
\midrule
\texttt{domain0} (OVC--shield gap) & warmPlate, OVC & fixed $T_{\mathrm{hot}}$ & grey diffuse, $\varepsilon=0.04$ \\
 & shield\_slave & fixed $T_{\mathrm{cold}}$ & grey diffuse, $\varepsilon=0.02$ \\
\texttt{domain1} (inner gap) & warmPlate & fixed $T_{\mathrm{hot}}$ & grey diffuse, $\varepsilon=0.04$ \\
 & coldPlate, shield\_bottom & fixed $T_{\mathrm{cold}}$ & grey diffuse, $\varepsilon=0.04$ \\
 & shield & fixed $T_{\mathrm{cold}}$ & grey diffuse, $\varepsilon=0.02$ \\
 & domain1\_to\_coax\_L$k$ & conjugate (mapped) & grey diffuse, $\varepsilon=0.02$ \\
\texttt{coax\_L1..L4} (solids) & warm cap / cold cap & fixed $T_{\mathrm{hot}}$ / $T_{\mathrm{cold}}$ & -- \\
 & lateral & conjugate to domain1 & -- \\
\bottomrule
\end{tabularx}\end{table}

\begin{table}[H]\centering\small
\caption{Gas-region settings that make the gaps behave as vacuum. Radiation in OpenFOAM's
\texttt{viewFactor} model is solved in fluid regions, so each gap is a fluid region made inert.}
\begin{tabular}{l l}
\toprule
Setting & Value \\
\midrule
pressure $p=p_{\mathrm{rgh}}$ & $10^{-3}$\,Pa (\texttt{pRefValue} $10^{-3}$) \\
gravity & $(0,0,0)$ \\
transport & constant, $\mu=10^{-11}$\,Pa\,s, $\Pr=0.67$ \\
density limits & $\rho_{\max}=10^{-4}$, $\rho_{\min}=10^{-12}$\,kg\,m$^{-3}$ \\
radiation solver frequency & every 10 iterations \\
agglomeration & \texttt{nFacesInCoarsestLevel} 250, \texttt{featureAngle} 20, \texttt{mergeLevels} 1 \\
\bottomrule
\end{tabular}\end{table}

\begin{table}[H]\centering\small
\caption{Harness representation. Cable UT-085-SS-SS: 304 stainless outer conductor
$A=1.585$\,mm$^2$, PTFE dielectric $A=2.001$\,mm$^2$, 304 stainless centre conductor
$A=0.205$\,mm$^2$. Material $k(T)$ from the NIST cryogenic fits
$\log_{10}k=\sum_n a_n(\log_{10}T)^n$. Each modelled cylinder carries $n_k$ lines with
$\kappa_{\mathrm{eff}}(T)=(n_k/A_{\mathrm{model}})\sum_iA_ik_i(T)$. $A_{\mathrm{model}}$ is the
conduction cross-section of the meshed cylinder in each stage span (calibrated from the mesh so
that $(A_{\mathrm{model}}/L)\int\kappa_{\mathrm{eff}}\,\mathrm dT$ equals the physical cable
integral $n_kL^{-1}\sum_iA_i\int k_i\,\mathrm dT$); it differs between the stage-1 and stage-2
geometries. Because $A_{\mathrm{model}}$ cancels in that product, the stage-integrated
conduction is fixed by the cable data by construction: the resolved model reproduces the
targets below to 0.01\,\%, which verifies the implementation but is not an independent
prediction.}
\setlength{\tabcolsep}{4pt}
\begin{tabular}{l r r r r r r}
\toprule
Span & $A_{\mathrm{model}}$ [m$^2$] & $L$ [m] & L1 drive, $n_k$=25 & L2 flux, 25 & L3 readout, 10 & L4 pump, 5 \\
\midrule
300$\to$50\,K & $2.975\times10^{-4}$ & 0.335 & 395.8 mW & 395.8 mW & 158.3 mW & 79.2 mW \\
300$\to$77\,K & $2.976\times10^{-4}$ & 0.335 & 370.0 mW & 370.0 mW & 148.0 mW & 74.0 mW \\
50$\to$4\,K & $3.934\times10^{-3}$ & 0.245 & 27.03 mW & 27.03 mW & 10.81 mW & 5.41 mW \\
77$\to$4\,K & $2.845\times10^{-3}$ & 0.245 & 62.33 mW & 62.33 mW & 24.93 mW & 12.47 mW \\
\bottomrule
\end{tabular}
\par\smallskip\footnotesize Per-bundle targets $n_kL^{-1}\sum_iA_i\int k_i\,\mathrm dT$; shares
38.5\,/\,38.5\,/\,15.4\,/\,7.7\,\% follow from $n_k$ with identical cable construction and length.
\end{table}

\begin{table}[H]\centering\scriptsize
\caption{Deployed $\kappa_{\mathrm{eff}}(T)=\sum_{n=0}^{5}c_nT^n$ [W\,m$^{-1}$K$^{-1}$, $T$ in K]
per bundle and stage span (\texttt{kappaCoeffs<8>}, last two coefficients zero). Fit error
$<0.5$\,\% over each span.}
\setlength{\tabcolsep}{4pt}
\begin{tabular}{l l r r r r r r}
\toprule
Span & Bundle & $c_0$ & $c_1$ & $c_2$ & $c_3$ & $c_4$ & $c_5$ \\
\midrule
300$\to$50\,K & L1 & $-2.1522\times10^{-1}$ & $3.1627\times10^{-2}$ & $-2.3184\times10^{-4}$ & $9.9921\times10^{-7}$ & $-2.1908\times10^{-9}$ & $1.9383\times10^{-12}$ \\
 & L2 & $-2.1519\times10^{-1}$ & $3.1622\times10^{-2}$ & $-2.3181\times10^{-4}$ & $9.9907\times10^{-7}$ & $-2.1905\times10^{-9}$ & $1.9380\times10^{-12}$ \\
 & L3 & $-8.6092\times10^{-2}$ & $1.2651\times10^{-2}$ & $-9.2741\times10^{-5}$ & $3.9970\times10^{-7}$ & $-8.7636\times10^{-10}$ & $7.7535\times10^{-13}$ \\
 & L4 & $-4.3044\times10^{-2}$ & $6.3253\times10^{-3}$ & $-4.6369\times10^{-5}$ & $1.9984\times10^{-7}$ & $-4.3816\times10^{-10}$ & $3.8766\times10^{-13}$ \\
300$\to$77\,K & L1 & $-1.3459\times10^{-1}$ & $2.9103\times10^{-2}$ & $-2.0205\times10^{-4}$ & $8.3171\times10^{-7}$ & $-1.7398\times10^{-9}$ & $1.4706\times10^{-12}$ \\
 & L2 & $-1.3457\times10^{-1}$ & $2.9099\times10^{-2}$ & $-2.0202\times10^{-4}$ & $8.3159\times10^{-7}$ & $-1.7396\times10^{-9}$ & $1.4704\times10^{-12}$ \\
 & L3 & $-5.3837\times10^{-2}$ & $1.1642\times10^{-2}$ & $-8.0824\times10^{-5}$ & $3.3269\times10^{-7}$ & $-6.9596\times10^{-10}$ & $5.8824\times10^{-13}$ \\
 & L4 & $-2.6917\times10^{-2}$ & $5.8206\times10^{-3}$ & $-4.0410\times10^{-5}$ & $1.6634\times10^{-7}$ & $-3.4796\times10^{-10}$ & $2.9411\times10^{-13}$ \\
50$\to$4\,K & L1 & $-4.8624\times10^{-4}$ & $9.1378\times10^{-4}$ & $3.6621\times10^{-5}$ & $-9.0343\times10^{-7}$ & $8.4798\times10^{-9}$ & $-2.8699\times10^{-11}$ \\
 & L2 & $-4.8624\times10^{-4}$ & $9.1378\times10^{-4}$ & $3.6621\times10^{-5}$ & $-9.0343\times10^{-7}$ & $8.4798\times10^{-9}$ & $-2.8699\times10^{-11}$ \\
 & L3 & $-1.9450\times10^{-4}$ & $3.6551\times10^{-4}$ & $1.4648\times10^{-5}$ & $-3.6137\times10^{-7}$ & $3.3919\times10^{-9}$ & $-1.1480\times10^{-11}$ \\
 & L4 & $-9.7247\times10^{-5}$ & $1.8275\times10^{-4}$ & $7.3242\times10^{-6}$ & $-1.8068\times10^{-7}$ & $1.6959\times10^{-9}$ & $-5.7398\times10^{-12}$ \\
77$\to$4\,K & L1 & $-7.4665\times10^{-4}$ & $1.2814\times10^{-3}$ & $4.9472\times10^{-5}$ & $-1.2245\times10^{-6}$ & $1.1714\times10^{-8}$ & $-4.2617\times10^{-11}$ \\
 & L2 & $-7.4665\times10^{-4}$ & $1.2814\times10^{-3}$ & $4.9472\times10^{-5}$ & $-1.2245\times10^{-6}$ & $1.1714\times10^{-8}$ & $-4.2617\times10^{-11}$ \\
 & L3 & $-2.9865\times10^{-4}$ & $5.1254\times10^{-4}$ & $1.9788\times10^{-5}$ & $-4.8978\times10^{-7}$ & $4.6855\times10^{-9}$ & $-1.7046\times10^{-11}$ \\
 & L4 & $-1.4932\times10^{-4}$ & $2.5626\times10^{-4}$ & $9.8938\times10^{-6}$ & $-2.4488\times10^{-7}$ & $2.3427\times10^{-9}$ & $-8.5230\times10^{-12}$ \\
\bottomrule
\end{tabular}\end{table}

\paragraph{Free-molecular gas conduction bound.} At $p<10^{-2}$\,Pa the mean free path
exceeds the gap and gas heat transport is free-molecular,
$\dot Q_{\mathrm{gas}}=\alpha\Lambda pA\Delta T$ \citep{37-Corruccini1959FreeMolecular}, with
$\Lambda=\frac{\gamma+1}{\gamma-1}\sqrt{R/(8\pi MT)}$ evaluated at the warm wall (helium,
$\gamma=5/3$: $\Lambda=2.10$\,W\,m$^{-2}$\,K$^{-1}$\,Pa$^{-1}$; nitrogen: 1.19),
accommodation coefficient $\alpha$, exchange area $A$ and $\Delta T=250$\,K. For the
enclosure's equivalent plate area $A=1.14$\,m$^2$ and helium with $\alpha=0.4$ the bound is
2.4\,mW at $10^{-5}$\,Pa, 24\,mW at $10^{-4}$\,Pa and 0.24\,W at $10^{-3}$\,Pa, that is 0.2,
2.3 and 23\,\% of the 1.03\,W harness conduction ($\alpha=1$: 0.6, 5.8 and 58\,\%). The model
therefore applies to a cryostat pumped to $10^{-4}$\,Pa or better, where the residual gas is a
budget line of a few per cent; the nominal $p=10^{-3}$\,Pa of the gap regions is a solver
setting for the radiation-only vacuum, not an operating pressure.

\begin{table}[H]\centering\footnotesize\setlength{\tabcolsep}{3pt}
\caption{View-factor agglomeration refinement and row normalisation. Closure = mean row sum
$\sum_jF_{ij}$ of the generated matrix (unity for a closed enclosure); loads with the matrix as
generated and, at the production level, with rows normalised.}
\begin{tabular}{l l r r r}
\toprule
Case & Coarsest-level faces & Closure (mean; range) & Radiation, as generated & Radiation, normalised \\
\midrule
300$\to$50\,K & 60 & -- & 9.46 W & -- \\
300$\to$50\,K & 250 (production) & 0.99--1.01; 0.50--1.13 & 8.41 W & 7.36 W \\
300$\to$50\,K & 400 & 0.98--0.99 & 6.49 W & 7.32 W \\
300$\to$50\,K & 800 & 0.90--0.96 & 3.26 W & 7.27 W \\
300$\to$50\,K & closed form & -- & \multicolumn{2}{c}{7.07 W} \\
300$\to$77\,K & 250 (production) & 0.99 & 8.38 W & 7.34 W (closed form 7.04) \\
50$\to$4\,K & 250 (production) & 1.00 & 5.29 mW & 3.16 mW (closed form 3.15) \\
77$\to$4\,K & 250 (production) & 1.00 & 29.7 mW & 17.8 mW (closed form 17.7) \\
\bottomrule
\end{tabular}\end{table}

\paragraph{Conservation and reciprocity audit.} Two residuals are reported for every view-factor
matrix. Global conservation: the net radiative power $\sum_{\mathrm{patches}}\int q_r\,\mathrm dA$
summed over all surfaces of a closed region, which must vanish; the generated matrices leave 1.5\,W
(400 faces) and 14\,W (800 faces) unbalanced on the 300$\to$50\,K case, the row-normalised ones
0.10, 0.13 and 0.36\,W at 250, 400 and 800 faces. Reciprocity: with the coarse-face areas $A_i$
rebuilt from the mesh by summing the fine faces of each agglomeration
(\texttt{scripts/reciprocity.py}), the area-weighted error $\sum|A_iF_{ij}-A_jF_{ji}|/\sum A_iF_{ij}$
of the normalised matrices is 3.9, 5.5 and 12.6\,\% at 250, 400 and 800 faces (outer gap region).
Row normalisation restores closure exactly but leaves reciprocity untouched, and both residuals grow
with agglomeration refinement, which is why the production level is the coarsest, 250 faces, and why
the 3--4\,\% departure from the closed form is quoted as the accuracy of the method.

\paragraph{Row normalisation.} Each fluid region of the enclosure model is a closed surface, so
every row of the view-factor matrix must sum to unity. The generated matrix (\texttt{viewFactorsGen},
Gauss quadrature tolerance 0.01) is stored sparsely, one entry per visible face pair; the
normalisation divides each row by its sum and leaves rows with a sum below 0.2 (none occurred)
untouched. The solve was then restarted from the converged production field for 600 further
iterations. Script \texttt{normalise\_F.py} in the repository.

\begin{table}[H]\centering\small
\caption{Emissivity sensitivity (production view factors) and MLI-equivalent variant; the
MLI-equivalent cases are also given with the row-normalised matrix, which is the value
carried into the main text and the ledger.}
\begin{tabular}{l r r r}
\toprule
Case & $\varepsilon$ warm / cold & Radiation [W] & Total [W] \\
\midrule
300$\to$50\,K, low & 0.02 / 0.01 & 5.45 & 6.48 \\
300$\to$50\,K, baseline & 0.04 / 0.02 & 8.41 & 9.44 \\
300$\to$50\,K, high & 0.06 / 0.03 & 11.48 & 12.51 \\
300$\to$77\,K, low & 0.02 / 0.01 & 5.44 & 6.40 \\
300$\to$77\,K, baseline & 0.04 / 0.02 & 8.38 & 9.35 \\
300$\to$77\,K, high & 0.06 / 0.03 & 11.44 & 12.40 \\
300$\to$50\,K, MLI-eq. & 0.04 / 0.003 & 1.81 & 2.84 \\
300$\to$77\,K, MLI-eq. & 0.04 / 0.003 & 1.81 & 2.77 \\
300$\to$50\,K, MLI-eq., rows normalised & 0.04 / 0.003 & 1.69 & 2.72 \\
300$\to$77\,K, MLI-eq., rows normalised & 0.04 / 0.003 & 1.68 & 2.65 \\
\bottomrule
\end{tabular}\end{table}

\section{Cold-plate sweep: full results}
\begin{landscape}
\begin{table}[H]\centering\small
\caption{All cold-plate runs at 9.4\,W (production mesh unless marked). $\dot m$ in g\,s$^{-1}$, volumetric flow in cm$^3$\,s$^{-1}$, $\Delta p$ in Pa, $\dot W_{\mathrm p}$ in $\mu$W, temperatures in K, $h$ in W\,m$^{-2}$K$^{-1}$; $T_{\mathrm{in}}$ = 45\,K (He), 77\,K (\LNtwo). $T_{\mathrm w}$ is the area mean over the heated-wall patch, $T_{\mathrm{w,max}}$ and $T_{\mathrm{w,min}}$ its extreme face values (\texttt{surfaceFieldValue} on the patch); margin = $T_{\mathrm{sat}}(5\,\mathrm{bar})-T_{\mathrm{w,max}}$ = $94.0\,\mathrm K-T_{\mathrm{w,max}}$; $y^+$ = maximum on the heated wall; closure = energy-balance error. Superscripts: g gravity on, c/f coarse/fine grid, v variable properties. Gravity-on nitrogen runs did not reach a steady state (see closure).}
\label{tab:S6}
\setlength{\tabcolsep}{3pt}\renewcommand{\arraystretch}{1.2}
{\begin{tabular}{l r r r r r r r r r r r r r r r r r r}
\toprule
Case (fluid, $p$, $\Re$) & $\dot m$ & $\dot V$ & $\Delta p$ & $\dot W_{\mathrm p}$ & $T_{\mathrm{out}}$ & $T_{\mathrm w}$ & $T_{\mathrm{w,max}}$ & $T_{\mathrm{w,min}}$ & $T_{\mathrm w}-T_{\mathrm b}$ & $h$ & $\Nu$ & $f$ & NTU & $\varepsilon$ & Gr/$\Re^2$ & margin & $y^+$ & closure \% \\
\midrule
He 1\,bar, 500 & 0.159 & 167 & 0.294 & 48.6 & 56.36 & 87.4 & 106.1 & 47.9 & 36.72 & 17.1 & 3.66 & 0.186 & 0.31 & 0.266 & 0.63 & -- & -- & -0.06 \\
He 1\,bar, 500$^{\mathrm{g}}$ & 0.159 & 167 & 0.331 & 54.6 & 56.21 & 67.5 & 95.2 & 47.9 & 16.87 & 37.1 & 7.96 & 0.21 & 0.674 & 0.49 & 0.29 & -- & -- & -1.41 \\
He 1\,bar, 500$^{\mathrm{v}}$ & 0.159 & 165 & 0.309 & 51.1 & 56.37 & 83.0 & 97.3 & 48.0 & 32.29 & 19.4 & 4.16 & 0.196 & 0.352 & 0.297 & 0.63 & -- & -- & -0.01 \\
He 1\,bar, 1000 & 0.318 & 316 & 0.671 & 211 & 50.68 & 76.8 & 92.6 & 47.1 & 28.94 & 21.7 & 4.64 & 0.112 & 0.196 & 0.178 & 0.15 & -- & -- & -0.05 \\
He 1\,bar, 1500 & 0.477 & 465 & 1.12 & 521 & 48.79 & 71.7 & 85.9 & 46.8 & 24.77 & 25.3 & 5.42 & 0.0847 & 0.153 & 0.142 & 0.06 & -- & -- & -0.05 \\
He 1\,bar, 2300 & 0.731 & 702 & 1.98 & 1\,388 & 47.47 & 67.1 & 79.7 & 46.6 & 20.82 & 30.1 & 6.45 & 0.0642 & 0.119 & 0.112 & 0.022 & -- & -- & -0.06 \\
He 1\,bar, 2300$^{\mathrm{g}}$ & 0.731 & 703 & 2.02 & 1\,419 & 47.48 & 62.7 & 80.9 & 46.6 & 16.41 & 38.2 & 8.18 & 0.0656 & 0.151 & 0.14 & 0.018 & -- & -- & 0.55 \\
He 1\,bar, 2300$^{\mathrm{c}}$ & 0.731 & 702 & 1.95 & 1\,367 & 47.47 & 67.1 & 79.9 & 47.3 & 20.85 & 30.1 & 6.44 & 0.0632 & 0.119 & 0.112 & 0.022 & -- & -- & -0.03 \\
He 1\,bar, 2300$^{\mathrm{f}}$ & 0.731 & 702 & 1.99 & 1\,400 & 47.47 & 67.0 & 79.7 & 46.2 & 20.79 & 30.1 & 6.46 & 0.0647 & 0.119 & 0.112 & 0.022 & -- & -- & -0.07 \\
He 1\,bar, 2300$^{\mathrm{v}}$ & 0.731 & 702 & 1.98 & 1\,390 & 47.47 & 66.0 & 76.9 & 46.7 & 19.80 & 31.7 & 6.78 & 0.0643 & 0.125 & 0.117 & 0.022 & -- & -- & -0.01 \\
He 1\,bar, 5000 & 1.59 & 1\,505 & 6.62 & 9\,955 & 46.14 & 55.3 & 61.3 & 46.5 & 9.73 & 64.4 & 13.8 & 0.0461 & 0.117 & 0.11 & 0.0023 & -- & 2.9 & -0.08 \\
He 1\,bar, 10000 & 3.18 & 2\,991 & 21.8 & 65\,113 & 45.57 & 50.3 & 54.3 & 46.4 & 5.00 & 125 & 26.8 & 0.0382 & 0.114 & 0.107 & 0.0003 & -- & 5.6 & -0.18 \\
He 5\,bar, 2300 & 0.731 & 140 & 0.395 & 55.5 & 47.47 & 67.1 & 79.8 & 46.6 & 20.82 & 30.1 & 6.45 & 0.0641 & 0.119 & 0.112 & 0.56 & -- & -- & -0.04 \\
He 18\,bar, 2300 & 0.731 & 39 & 0.11 & 4.29 & 47.47 & 67.1 & 79.7 & 46.6 & 20.82 & 30.1 & 6.45 & 0.0643 & 0.119 & 0.112 & 7.3 & -- & -- & -0.04 \\
\LNtwo 3\,bar, 500 & 4.09 & 5.06 & 0.202 & 1.02 & 78.13 & 85.0 & 89.2 & 77.7 & 7.46 & 84 & 5.77 & 0.163 & 0.151 & 0.14 & 40 & 4.8 & -- & 0.00 \\
\LNtwo 3\,bar, 500$^{\mathrm{g}}$ & 4.09 & 5.06 & 1.68 & 8.5 & 77.32 & 78.9 & 81.7 & 77.7 & 1.74 & 361 & 24.8 & 1.35 & 0.65 & 0.478 & 9.4 & 12.3 & -- & -71.96 \\
\LNtwo 3\,bar, 1000 & 8.17 & 10.1 & 0.501 & 5.07 & 77.56 & 83.1 & 86.6 & 77.6 & 5.77 & 109 & 7.45 & 0.101 & 0.0977 & 0.0931 & 7.8 & 7.4 & -- & 0.00 \\
\LNtwo 3\,bar, 1500 & 12.3 & 15.2 & 0.869 & 13.2 & 77.38 & 82.1 & 85.3 & 77.5 & 4.92 & 127 & 8.75 & 0.0778 & 0.0765 & 0.0736 & 2.9 & 8.7 & -- & 0.00 \\
\LNtwo 3\,bar, 2300 & 18.8 & 23.3 & 1.57 & 36.6 & 77.25 & 81.2 & 84.1 & 77.5 & 4.12 & 152 & 10.4 & 0.0599 & 0.0595 & 0.0578 & 1 & 9.9 & -- & -0.00 \\
\LNtwo 3\,bar, 2300$^{\mathrm{g}}$ & 18.8 & 23.3 & 3 & 69.8 & 77.22 & 78.5 & 79.9 & 77.5 & 1.35 & 463 & 31.8 & 0.114 & 0.181 & 0.166 & 0.34 & 14.1 & -- & -10.77 \\
\LNtwo 3\,bar, 2300$^{\mathrm{c}}$ & 18.8 & 23.3 & 1.55 & 36.1 & 77.25 & 81.2 & 84.1 & 77.7 & 4.09 & 153 & 10.5 & 0.0592 & 0.06 & 0.0582 & 1 & 9.9 & -- & 0.01 \\
\LNtwo 3\,bar, 2300$^{\mathrm{f}}$ & 18.8 & 23.3 & 1.58 & 36.8 & 77.25 & 81.3 & 84.1 & 77.4 & 4.14 & 151 & 10.4 & 0.0603 & 0.0592 & 0.0575 & 1.1 & 9.9 & -- & -0.02 \\
\LNtwo 3\,bar, 5000 & 40.9 & 50.6 & 5.57 & 282 & 77.11 & 78.6 & 80.0 & 77.4 & 1.55 & 405 & 27.8 & 0.0449 & 0.073 & 0.0704 & 0.083 & 14.0 & 2.9 & 0.04 \\
\LNtwo 3\,bar, 10000 & 81.7 & 101 & 18.6 & 1\,878 & 77.06 & 77.8 & 78.6 & 77.4 & 0.81 & 777 & 53.3 & 0.0374 & 0.0699 & 0.0675 & 0.011 & 15.4 & 5.6 & 0.03 \\
\bottomrule
\end{tabular}}
\end{table}

\end{landscape}

\begin{table}[H]\centering\small
\caption{Grid convergence index study for both fluids at $\Re=2300$ (Celik et al.\ 2008); meshes 86\,400 /
291\,600 / 984\,150 cells, refinement ratio 1.5, production mesh = medium.}
\begin{tabular}{l l r r r r r r r}
\toprule
Fluid & Quantity & Coarse & Medium & Fine & $p$ & Richardson & GCI$_{\mathrm{medium}}$ & GCI$_{\mathrm{fine}}$ \\
\midrule
He, 1\,bar & $\Nu$ & 6.439 & 6.449 & 6.457 & 0.64 & 6.485 & 0.70\,\% & 0.54\,\% \\
He, 1\,bar & $f$ & 0.0632 & 0.0642 & 0.0647 & 1.40 & 0.0655 & 2.5\,\% & 1.4\,\% \\
He, 1\,bar & $h$ [W\,m$^{-2}$K$^{-1}$] & 30.05 & 30.10 & 30.14 & 0.64 & 30.27 & 0.70\,\% & 0.54\,\% \\
\LNtwo, 3\,bar & $\Nu$ & 10.513 & 10.436 & 10.382 & 0.87 & 10.26 & 2.2\,\% & 1.5\,\% \\
\LNtwo, 3\,bar & $f$ & 0.0592 & 0.0599 & 0.0603 & 1.71 & 0.0607 & 1.6\,\% & 0.8\,\% \\
\LNtwo, 3\,bar & $h$ [W\,m$^{-2}$K$^{-1}$] & 153.2 & 152.1 & 151.3 & 0.87 & 149.5 & 2.2\,\% & 1.5\,\% \\
\bottomrule
\end{tabular}\end{table}

\paragraph{Property correction and energy accounting.} The helium transport properties of the
earlier version of this comparison ($\mu=2.73\,\mu$Pa\,s, $k=0.0185$\,W\,m$^{-1}$K$^{-1}$)
correspond to helium near 15\,K; the present runs use the 50\,K values ($6.36\,\mu$Pa\,s,
0.0467\,W\,m$^{-1}$K$^{-1}$, CoolProp/NIST). The mixing-cup outlet temperature is the
mass-flux-weighted average $\sum\phi_fT_f/\sum\phi_f$ over the outlet faces
(\texttt{surfaceFieldValue}, \texttt{weightedAverage} with \texttt{phi}); the
\texttt{weightedAreaAverage} operation, which multiplies the mass flux by the face area a second
time, biases the outlet temperature low on a graded mesh and was the cause of the 2--10\,\% energy
imbalances reported earlier. The delivered wall heat flux was verified with the solver-native
\texttt{wallHeatFlux} function object (9.395\,W).

\paragraph{Wall temperatures and the nitrogen loop pressure.} Wall temperatures are read from
the heated-wall patch with \texttt{surfaceFieldValue} (\texttt{max}, \texttt{min},
\texttt{areaAverage}; include file \texttt{system/wallTstats} in every case), and the
wall-profile figure of the manuscript uses the patch face values along the centre row of the
first channel (\texttt{postProcessing/wallFaces/wall\_centreline.csv}). A line sample 0.4\,mm
inside the fluid, which an earlier post-processing used, underestimates the hottest wall face
by 8--10\,K for helium and 1--3.5\,K for nitrogen and is not used anywhere in the paper. The
nitrogen cases are stored under the name \texttt{LN2\_3bar\_*} because their properties were
generated at 3\,bar; the Boussinesq fields do not depend on the loop pressure (density,
viscosity and conductivity of the 77\,K liquid change by 0.06, 0.3 and 0.1\,\% between 3 and
5\,bar), and the saturation margins of Table~S7 are evaluated at the 5\,bar loop pressure of
the manuscript, with the 3\,bar value (87.9\,K) quoted where it matters.

\section{Hardware efficiencies used in the real ledger}
\begin{table}[H]\centering\small
\caption{Stage efficiencies and the nitrogen supply-energy multiplier used in the real ledger.}
\begin{tabularx}{\textwidth}{l l X}
\toprule
Stage & $\eta$ range & Source \\
\midrule
4\,K cryocooler stage & 0.008--0.015 & PT415: 1.5\,W at 4.2\,K with 40\,W at 45\,K for 10.7\,kW
\citep{Bluefors2026PT415}: the two stages together return 0.031 of their combined Carnot work, which the
per-stage ranges reproduce to within 7\,\% at their midpoints;
Radebaugh: ``at 4\,K efficiencies of about 1\,\% of Carnot or less are typical'' \citep{Radebaugh2009Cryocoolers} \\
He first stage (45--80\,K) & 0.10--0.15 & Radebaugh: 10--15\,\% of Carnot for the best commercial 80\,K
coolers; large plants up to 30\,\% \citep{Radebaugh2009Cryocoolers, Claudet2010LHCExergy} \\
LN$_2$ supply (basis of Table 6) & $M_{\mathrm{N_2}}$ = 9.1--12.7 W per W intercepted & plant
consumption 0.5--0.7\,kWh\,kg$^{-1}$ \citep{Karabuga2018N2Liquefaction} divided by the latent
heat 199\,kJ\,kg$^{-1}$ at 1\,bar; equivalent to 0.30--0.40 of Carnot at 77\,K (ideal
liquefaction work 0.213\,kWh\,kg$^{-1}$ \citep{BarronNellis2016Cryogenic}); the
fraction-of-Carnot form gives 7.2--9.7\,W/W and changes the nitrogen totals by 1--2\,\% \\
\bottomrule
\end{tabularx}\end{table}

\bibliography{precooling_refs}
\end{document}
